%% 03-porter-compiler.tex --- included by ../sfc.tex
\section{The \PORTER{} Compiler}
\label{sec:sfc}

Given a circuit and an initial layout $\ell_0$ (which \PORTER{} builds
itself, Section~\ref{sec:sfc-layout}), \PORTER{} repeats the following
loop (Algorithm~\ref{alg:sfc}) until every gate has been emitted or a
configurable iteration budget is exhausted. Algorithm~\ref{alg:sfc}
charts one iteration; its shaded region is the plan-first remote stage
that distinguishes \PORTER{} from compilers deciding at
single-action granularity; \appExample{} walks through it on a concrete
front where it changes the outcome.


\begin{algorithm}[t]
\caption{\PORTER{} \texttt{compile}$(C, \ell_0)$}
\label{alg:sfc}
\begin{algorithmic}[1]
\STATE $\ell \leftarrow \ell_0$; initialise DAG active/frontier sets from $C$
\WHILE{DAG has active gates}
  \STATE $F \leftarrow$ DAG front (bounded by \texttt{dag\_front\_limit})
  \STATE classify $F$ into singles, direct-local CX, local-CX-needing-SWAP,
         adjacent-remote CX, non-adjacent-remote CX \hfill\textit{(Sec.~\ref{sec:sfc-classify})}
  \IF{singles or direct-local CX present}
    \STATE execute them all; remove from DAG; \textbf{continue}
  \ENDIF
  \IF{local-CX-needing-SWAP present}
    \STATE generate candidate intra-core SWAPs, apply the best-scoring one
           (Eq.~\ref{eq:sfc-local}); \textbf{continue} \hfill\textit{(Sec.~\ref{sec:sfc-intra})}
  \ENDIF
  \STATE generate normal remote plans $\mathcal{P}$ and blocked-teleport
         plans $\mathcal{B}$ \hfill\textit{(Sec.~\ref{sec:sfc-scoring})}
  \IF{$\mathcal{P}\neq\emptyset$}
    \STATE $p^\ast \leftarrow$ best ready-or-preparable plan in $\mathcal{P}$
    \IF{some $b\in\mathcal{B}$ has gain $\gg$ gain$(p^\ast)$}
      \STATE $p^\ast \leftarrow$ best opportunistic capacity-relief plan for $b$
             \hfill\textit{(Sec.~\ref{sec:sfc-relief})}
    \ENDIF
  \ELSE
    \STATE BFS from each blocked target core to its nearest reservoir;
           $p^\ast \leftarrow$ best tail-relief plan \hfill\textit{(Sec.~\ref{sec:sfc-relief})}
  \ENDIF
  \IF{$p^\ast$ is \texttt{None}} \STATE \textbf{report deadlock; abort} \ENDIF
  \STATE apply $p^\ast$'s own preparation SWAPs until ready, then execute it
\ENDWHILE
\RETURN emitted operations, final layout, stats
\end{algorithmic}
\end{algorithm}

\subsection{Initial Layout}
\label{sec:sfc-layout}

\PORTER{} builds its starting layout $\ell_0$ with one deterministic,
seed-free heuristic --- no external partitioner, no best-of-$N$ search
--- and respects the capacity invariant the routing loop depends on:
every core keeps at least $f_\mathrm{min}{=}2$ free positions, so from
the first iteration any core can serve as a \TeleData{} destination and
as a relief reservoir.

\textbf{Temporal locality-greedy.} Three deterministic steps.
\emph{(i) Weighting}: each CX gate $g=(q_i,q_j)$ contributes
$\gamma^{\ell(g)}$ to the interaction weight $w(q_i,q_j)$, where
$\ell(g)$ is $g$'s as-soon-as-possible layer in the CX-only DAG and
$\gamma{=}0.99$, so early interactions weigh more than late ones while
repeated interactions accumulate. \emph{(ii) Core assignment}: logical
qubits are visited in the order they first appear as a CX operand and
each is placed in the core holding its heaviest already-placed
neighbour when that core has headroom above $f_\mathrm{min}$, falling
back to the core with the most remaining reserved capacity.
\emph{(iii) Physical placement}: within each core, logicals are laid
onto positions by descending interaction weight onto ascending
intra-core centrality, so the most-connected qubits take the most
central slots.

The steps are complementary. The temporal discount front-loads the
opening layers, co-locating the pairs that must interact soonest;
accumulated multiplicity plus first-appearance order tracks
chain-structured circuits, where only one disjoint pair is exposed per
layer and a balanced-first-layer rule would scatter the chain across
cores; and the centrality step keeps hub qubits cheap to reach. The
construction costs microseconds and consumes no randomness, so
\PORTER{} has no seed to report.

\subsection{Backward Refinement}
\label{sec:sfc-refine}

\PORTER{}'s default protocol is one forward pass from $\ell_0$. We also
report an optional refinement schedule, because
Section~\ref{sec:shared_results} shows a fixed layout is not the whole
difference between the two compilers: given the \emph{same} externally
supplied layout, \dSABRE{}'s forward--backward--forward schedule
extracts more from it than a single forward pass does.

The schedule is \SABRE{}'s bidirectional layout
refinement~\cite{li2019tackling} applied to whole routes. Let $R(C,\ell)$
denote routing $C$ from $\ell$ and $\ell'$ the layout it ends on.
Pass~1 computes $R(C,\ell_0)$; pass~2 computes
$R(C^{\mathrm{rev}},\ell_1')$ on the reversed circuit and
\emph{discards its route}, keeping only $\ell_2'$; pass~3 computes
$R(C,\ell_2')$. Reversing $C$ reverses every dependency edge, so
routing it from the forward pass's end state yields a layout adapted to
the \emph{start} of the forward circuit. Odd pass counts above three
repeat the cycle; $P{=}5$ throughout. We report the best forward pass
by EPR, so refinement can never report worse than the single pass. It
adds no randomness and its whole cost is routing passes --- $P$ calls
where the default makes one, against \dSABRE{}'s nine.

It recovers less here than it does for \dSABRE{}, for a structural
reason: refinement consumes a post-routing layout, and routing destroys
the free-slot reservation $\ell_0$ was built with, so intermediate
passes often regress and only the best-pass rule keeps that harmless.
Where it pays it pays well --- 64q~QNN $864 \to 537$, 64q~Multiplier
$1{,}358 \to 763$.

\subsection{Front Classification}
\label{sec:sfc-classify}

Every front gate is classified under the current layout $\ell$. 1Q
gates are always directly executable; a CX $g=(q_i,q_j)$ with both
operands in core $c$ is \emph{direct-local} if
$\dL(c,\ell(q_i),\ell(q_j))=1$ and \emph{local-needing-SWAP}
otherwise; a CX whose operands lie in different cores is
\emph{adjacent-remote} if those cores are core-graph neighbours and
\emph{non-adjacent-remote} otherwise. Directly executable gates drain
first and completely, intra-core SWAPs are exhausted next, and
remote-stage decisions are taken only once no front gate is
intra-core. The front is size-capped at $96$ gates to bound
per-iteration cost, a cap that never binds here since front gates are
pairwise qubit-disjoint.

\subsection{Intra-Core SWAP Scoring}
\label{sec:sfc-intra}

For a local-needing-SWAP gate set $F_\mathrm{loc}$, \PORTER{} builds
candidate SWAP edges from the coupling-graph neighbours of every
physical position touched by $F_\mathrm{loc}$, simulates each SWAP,
and scores it with \SABRE{}'s front-plus-lookahead
form~\cite{li2019tackling}:
\begin{equation}
  J(a) = \bar d_{\mathcal{F}}(a) + \lambda_\mathrm{loc}\,\bar d_{\mathcal{E}}(a),
  \qquad \lambda_\mathrm{loc}=0.8,
  \label{eq:sfc-local}
\end{equation}
where $\bar d_{\mathcal{F}}(a)$ is the mean, over the gates
$(q_i,q_j)\in F_\mathrm{loc}$, of the intra-core distance
$\dL\bigl(c,\ell_a(q_i),\ell_a(q_j)\bigr)$ under the post-SWAP layout
$\ell_a$, and $\bar d_{\mathcal{E}}(a)$ is the same mean taken over an
extended set $\mathcal{E}_\mathrm{loc}$ of up to
$N_\mathrm{loc}{=}20$ local CX gates collected from the next
$L_\mathrm{loc}{=}12$ DAG layers via a virtual frontier simulation
(Section~\ref{sec:sfc-simstate}). Ties are broken first by the raw
front-layer distance, then, if still tied, by a decision uniformly
sampled from a constant-seeded random-number generator, with every
such tie-break recorded in \texttt{RouterStats} for reproducibility
auditing.

\subsection{Plan-First Remote Stage}
\label{sec:sfc-plan-first}

\PORTER{} manipulates remote work as \emph{plans}. A plan is a
complete, committable specification of one inter-core action: which
primitive fires, which gate or same-control run it resolves, its
direction, and consequently its \emph{readiness conditions} --- which
logical endpoint must sit adjacent to which communication port, and
which ports must be free. What it deliberately leaves unspecified is
the SWAP sequence that establishes those conditions. Alternatives are
compared at the level of \emph{what} remote work to do next, each
scored as if its preconditions were met; only the committed plan is
handed to the machinery that works out \emph{how}.

Both baselines decide at a finer grain. \TeleSABRE{}'s pool holds
staging SWAPs and remote operations together, so the choice is made
implicitly, one micro-move at a time: every staging SWAP it selects is
also a vote for whichever remote action that SWAP happens to serve,
votes for competing actions interleave, and staging effort is routinely
spent on actions never taken. \dSABRE{} does separate the stages, but
its candidate is a single teleport hop rather than a plan, and its
score adds that hop's staging and eviction cost directly to the
progress it buys, so a hop cheap to stage now can outrank one that
resolves more gates. \PORTER{} ranks by gain alone, demotes readiness
cost to a tie-break, and refuses to reopen the choice until the
selected plan is ready (\appExample{}); Section~\ref{sec:ablation}
measures what that persistence is worth.

When the front holds only remote CX gates, \PORTER{} generates two
disjoint candidate sets. \textbf{Normal plans} $\mathcal{P}$: for every
adjacent-remote CX, a \TeleGate{}; a \TeleGatePacket{} if a
same-control run of $\ge 2$ remote CX gates can be collected by
forward DAG simulation (Section~\ref{sec:sfc-packet}); and, for either
endpoint, a \TeleData{} plan into the partner's core if that core has
at least $f_\mathrm{min}{=}2$ free positions. For non-adjacent-remote
CX gates, only distance-reducing \TeleData{} hops with sufficient
destination capacity. \textbf{Blocked plans} $\mathcal{B}$: the same
\TeleData{} candidates for destinations that currently fail the
$f_\mathrm{min}$ test. These are never executed; they exist to detect
when no normal plan exists because every destination is saturated, or
when one is available but a blocked plan promises much more if capacity
were relieved first.

Every plan is scored (Section~\ref{sec:sfc-scoring}) \emph{before} any
preparation SWAP for it exists. \PORTER{} then selects the best plan in
$\mathcal{P}$ that is ready or admits an improving preparation SWAP;
compares its gain against the best of $\mathcal{B}$ and, if the blocked
plan leads by more than $m_\mathrm{assist}{=}5.0$, runs opportunistic
capacity relief for that saturated core instead
(Section~\ref{sec:sfc-relief}); and, if $\mathcal{P}$ is empty, runs
the same breadth-first relief search directly on $\mathcal{B}$'s
destinations, shortest hop count first. Once committed, it generates
and applies \emph{only that plan's} preparation SWAPs until the plan is
ready, no improving SWAP remains, or the full decision state --- plan
identity, layout, RNG state --- repeats, which raises rather than
looping.

\textbf{Preparation SWAPs.}\label{sec:sfc-prep} A committed plan's
readiness requirements are exactly two \emph{(core, neighbour-core,
endpoint)} triples (source and destination side for \TeleData{};
control and target side for \TeleGate{}/\TeleGatePacket{}, whose
target-side routing is deferred to execution time,
Section~\ref{sec:sfc-packet}). For each unmet requirement, candidate
preparation SWAPs are restricted to edges incident to either the
plan's logical endpoint or a nearest free physical position tied for
minimum distance to the relevant communication port (only while that
port remains occupied); each candidate is scored by the post-SWAP
value of the same readiness-cost function used to rank blocked plans
(Section~\ref{sec:sfc-scoring}), and at most one improving SWAP per
core side is applied per pass.

\subsection{Plan Scoring: Direct, Latent, and Lookahead Gain}
\label{sec:sfc-scoring}

Every plan scores $\mathrm{total} = \mathrm{direct} +
\mathrm{lookahead}$, computed per plan kind but sharing one
$\Delta$-form convention: for a CX gate $g=(q_i,q_j)$ and a candidate
relocation, $\Delta(g) = \dC(\kappa(q_i),\kappa(q_j)) -
\dC(\kappa'(q_i),\kappa'(q_j))$, the reduction in core-graph distance
between $g$'s two cores.

A \textbf{\TeleGate{}} contributes a constant direct gain of $1$ --- one
remote CX resolved --- and no lookahead. A \textbf{\TeleGatePacket{}}
contributes a direct gain equal to the number of remote CX gates in the
collected packet; every gate in an accepted packet is current rather
than speculative work, since the whole packet executes under one
cat-entangler resource once started. A \textbf{\TeleData{}} plan
decomposes into three tiers, all computed by simulating the DAG
frontier under the hypothetical post-teleport core assignment:
\emph{direct} --- front-layer CX gates already involving the moving
qubit, plus the seed gate; \emph{latent-direct} --- CX gates on the
moving qubit that surface while the simulator can still progress using
only 1Q and now-local gates, which are current pressure rather than
prediction because the candidate teleport fixes their core distance the
moment they appear; and \emph{lookahead} --- gates on the moving qubit
found in up to $L_\mathrm{tp}$ further simulated layers, a gate in layer
$k$ entering with weight
$w_k = \max(w_\mathrm{min}, w_0\gamma^{k-1})$, $w_0{=}0.9$,
$\gamma{=}0.9$, $w_\mathrm{min}{=}0.5$. Thus
\begin{align}
  \mathrm{direct} &= \!\!\sum_{g\in\mathrm{direct}\cup\mathrm{latent}}\!\!\Delta(g),
  &
  \mathrm{lookahead} &= \!\!\sum_{(g,k)\in\mathrm{la}}\!\!w_k\,\Delta(g),
  \label{eq:sfc-teleport}
\end{align}
with latent-direct gates entering at full weight alongside direct ones.
A capacity-risk term $1/\max(1,f_\mathrm{dst}{-}1)$ is tracked
alongside the score for tie-breaking, but never added to it.

This is the point at which \PORTER{} differs most sharply from
\dSABRE{}: preparation cost appears nowhere in $\mathrm{total}$. Plans
are ranked on gain, and readiness cost enters only as the third
tie-break in the selection key, after gain and plan kind. The horizon $L_\mathrm{tp}$ is counted in
\emph{layers}, and the collection admits at most one gate on the moving
qubit per layer, because a qubit has at most one zero-predecessor gate at
any time. A gate cap $N_\mathrm{tp}$ therefore cannot bind whenever
$N_\mathrm{tp}\ge L_\mathrm{tp}$, and the horizon is the single
parameter that governs how far the estimate looks. No one value of
$L_\mathrm{tp}$ serves every circuit --- Section~\ref{sec:horizon}
shows the best value spans $4$ to $128$ across the benchmark --- so
\PORTER{} selects it per circuit rather than fixing it.

\subsection{Same-Control Packet Collection}
\label{sec:sfc-packet}

Same-control CX gates are chained by DAG dependency edges, so they
never appear in the front together: a packet cannot be read off the
front and must be collected by forward simulation.
From a seed adjacent-remote CX $g=(q_c,q_t)$ with control core $c_c$
and target core $c_t$, \PORTER{} simulates the DAG frontier forward for
up to $120$ layers, greedily absorbing any front CX sharing control
$q_c$ whose target is also in $c_t$ (up to $8$ such gates) and any
front 1Q gate whose qubit is already in $c_t$, since those execute
freely while the cat resource is live. Collection stops at the first
unproductive layer or at the size cap. The window counts simulated \emph{layers}, not collected CXs, because
same-control CXs are DAG-chained and surface one at a time; in practice
it is a cost guard, not an active knob. Because the simulation admits gates from later
generations than the bounded current front, a packet may contain gates
a naively front-limited scan could not reach. The collected order is
replayed exactly at commit time as one \texttt{cat\_entangler}
(consuming the packet's single pair), one \texttt{cat\_cx} per member
--- each target routed to its port by ordinary local SWAPs forbidden
from touching the two live ports --- and a closing
\texttt{cat\_disentangler}.

\subsection{Capacity Relief}
\label{sec:sfc-relief}

\textbf{Tail-relief BFS.} When every remote plan is capacity-blocked,
\PORTER{} runs a breadth-first search from each blocked destination core
to the nearest \emph{reservoir} core --- one with at least $f_\mathrm{min}$
free slots --- optionally chained beyond one hop
(\texttt{capacity\_relief\_chain\_enabled}). Every logical qubit
currently in the BFS path's penultimate core becomes a tail-relief
\TeleData{} candidate into the reservoir, scored and selected by the
same plan machinery, grouped and tried in ascending hop-count order.

\textbf{Opportunistic relief.} Even when a normal plan is available,
a severely blocked plan can still win if relieving its target core's
capacity would unlock substantially more gain
(Section~\ref{sec:sfc-plan-first}); this differs from tail-relief BFS
in firing pre-emptively, before any plan is actually starved, whenever
the margin $m_\mathrm{assist}$ is exceeded.

\subsection{Progress and the Failure Model}
\label{sec:sfc-progress}

Deterministic greedy routing on a capacity-constrained machine can
stall, and the mechanisms above are best read as one answer to that
risk. \PORTER{} layers four defences.

\emph{(i) Keep capacity from reaching zero.} A core with no free
physical qubit is terminal under the port-based action model: it can
never vacate an occupied communication port with local SWAPs, so no
inter-core action can touch it and its qubits are stranded. \PORTER{}
legalises the entry layout so no core starts full, and admits a
\TeleData{} plan into core $c$ as a \emph{normal} plan only while
$\mathrm{free}(c) \ge f_\mathrm{min}{=}2$. Ordinary routing therefore
cannot drive a core into the terminal state; only an externally
supplied layout can present one, which is why \appPytket{}'s
pytket-dqc source is the only place the compiler reports
\emph{infeasibility}.

\emph{(ii) Detect starvation rather than thrash.} An empty plan set
alongside a non-empty $\mathcal{B}$ is an explicit signal that the
front is starved \emph{by capacity}, and it names the cores
responsible. \emph{(iii) Manufacture capacity where it is missing.}
Tail-relief BFS turns that signal into work
(Section~\ref{sec:sfc-relief}). \emph{(iv) Bound every loop.} Plan
sides whose preparation score is infinite are never selected; during
preparation the compiler hashes its full decision state once per round
and raises if it repeats; two budgets bound the loops from outside; and
when no relief path reaches a reservoir \PORTER{} raises a diagnostic
naming the saturated cores.

We state the resulting guarantee precisely, because it is weaker than
it may appear. \PORTER{} is \emph{safe} but not \emph{complete}: it
never emits an invalid or partial circuit and always terminates, but
termination may be by exception, and it has no analogue of
\dSABRE{}'s termination-with-success proof
(Section~\ref{sec:related}). Its completeness is empirical, and
Section~\ref{sec:ablation} measures what buys it: with both relief
tiers active none of the twenty-four instances deadlocks in any
protocol; with both disabled, four of the six 80-qubit instances do.
No instance in this evaluation is left unrouted, but that is a
property of the protocol rather than of the router, and the exception
is instructive: at a fixed lookahead horizon, Multiplier on the
heavy-hex star of Section~\ref{sec:heavyhex} livelocks in the
intra-core SWAP loop --- capacity never reaches zero, no plan is
starved, the diagnostic never fires, and it is the iteration budget
that ends the run. Routing a second horizon escapes it
(Section~\ref{sec:horizon}), which is a way of not encountering the
failure rather than of ruling it out. Supplying a comparable
worst-case guarantee is the clearest piece of future work this design
leaves open.

A per-iteration cost analysis of the loop is given in \appComplexity{}.

%%────────────────────────────────────────────────────────────────────────────
